\subsection{Source and physical adaptation}\label{sec:absolute-quality}
We train paired physical updates for both source families using the same eight training compositions and matched maximum proposal and training budgets. Table~\ref{tab:absolute-quality} compares the four models on 24 compositions; Appendix~\ref{sec:factorial-details} details preparation costs and evaluation reuse.

\begin{table}[t]
\centering
\caption{Source and paired-update comparison under GFN2, with 1,536 raw attempts per model. Force medians use graph-valid, successfully scored outputs. Joint yield requires graph validity and force RMS at most 5 eV/\AA, with all attempts in the denominator.}
\small
\begin{tabular}{lrrr}
\toprule
Model & Graph validity (\%) & Median force (eV/\AA) & Joint yield (\%)\\
\midrule
Gaussian source &44.53&3.404&35.68\\
Gaussian + paired update &46.61&3.222&40.63\\
Harmonic source &52.73&3.287&44.53\\
Harmonic + paired update &53.45&3.037&47.59\\
\bottomrule
\end{tabular}\label{tab:absolute-quality}
\end{table}

GFN2 joint yield increases by 4.95 points for Gaussian (composition 95\% interval [2.86,7.10]) and 3.06 for harmonic ([1.04,5.14]), with both runs improving for each source. On jointly valid pairs, Gaussian adaptation also lowers energy by 0.01465 eV/atom and force RMS by 0.4618 eV/\AA.

After adaptation, harmonic exceeds Gaussian by 6.84 graph-validity points ([4.23,9.51]) and 6.97 joint-yield points ([3.26,10.61]). Graph validity improves in both runs, but joint-yield differences are +16.28 and $-2.34$, showing substantial run variation. eSEN likewise favors adaptation within each source and harmonic in the pooled comparison. Figure~\ref{fig:primary-quality} in Appendix~\ref{app:components} shows the threshold dependence.
